\documentclass[conference]{IEEEtran}
\IEEEoverridecommandlockouts
% The preceding line is only needed to identify funding in the first footnote. If that is unneeded, please comment it out.
\usepackage{cite}
\usepackage{amsmath,amssymb,amsfonts}
\usepackage{algorithmic}
\usepackage{graphicx}
\usepackage{textcomp}
\usepackage{xcolor}
\usepackage{url}
\usepackage{listings}
\usepackage{graphicx} % Required for \resizebox
\usepackage{array} % Required for column width adjustments
\usepackage{tabularx}
\usepackage{balance}
\usepackage{colortbl} 
\usepackage{booktabs}
\usepackage{dialogue}
\usepackage{multirow} 
% \usepackage{CJK}
\usepackage{CJKutf8}
\usepackage[utf8]{inputenc}
\usepackage{minted}
\usepackage{caption}
\usepackage{CJK}
\usepackage{longtable}
\usepackage{array}


\def\BibTeX{{\rm B\kern-.05em{\sc i\kern-.025em b}\kern-.08em
    T\kern-.1667em\lower.7ex\hbox{E}\kern-.125emX}}
\begin{document}

\title{Enhancing Logical Reasoning in LLMs through Few-Shot Prompting and Fine-Tuning: A Case Study with Turtle Soup Puzzles}
%\title{Enhancing Logical Reasoning in LLMs through Few-Shot Prompting and Fine-Tuning
%%: A Case Study with Turtle Soup Puzzles
%}
% https://challenge.ai.mgtv.com/#/track/25

\author{Anonymous Submission}

\iffalse
\author{\IEEEauthorblockN{Donghao HUANG\textsuperscript{1,2}, Zhaoxia WANG\textsuperscript{1}
}

\IEEEauthorblockA{School of Computing and Information Systems, Singapore Management University, Singapore\textsuperscript{1}}
\IEEEauthorblockA{Research and Development, Mastercard, Singapore\textsuperscript{2}}
\IEEEauthorblockA{dh.huang.2023@smu.edu.sg; zxwang@smu.edu.sg}
%\thanks{\IEEEauthorrefmark{1}
% Corresponding Authors: Zhaoxia WANG (e-mail: zxwang@smu.edu.sg)
}
\fi

\maketitle

\begin{abstract}

Recent advancements in Large Language Models (LLMs) have greatly improved their ability to tackle complex reasoning tasks. This study evaluates the performance of several leading LLMs, including OpenAI’s GPT-4o, Meta’s Llama3.1, Mistral AI's Mistral-7B, Alibaba’s Qwen2 and Qwen2.5, and Shanghai AI Laboratory’s InternLM2.5, on logical reasoning challenges presented by the Turtle Soup game—a lateral thinking puzzle that demands deductive reasoning. We conducted extensive experiments using both few-shot prompting and fine-tuning, and introduced the Valid Classification Ratio (VCR) as a novel metric to assess the reliability of model outputs. Our findings reveal that while few-shot learning can be effective, fine-tuning generally yields superior results. Qwen2.5-72B achieved the highest F1 score of 80.88\% through few-shot learning, surpassing GPT-4o's best score of 80.36\%. Llama3.1-70B also outperformed GPT-4o, with an F1 score of 80.77\% after fine-tuning. Additionally, InternLM2.5-7B-1M, featuring an extended context length of 1 million tokens, demonstrated significant improvement, achieving 77.01\% in few-shot and 80.28\% after fine-tuning, ranking third overall among fine-tuned models. These results underscore the substantial potential of latest open-source models in complex reasoning tasks, showing that they can not only match but also surpass proprietary alternatives.

% Recent advancements in Large Language Models (LLMs) have significantly enhanced their ability to tackle complex reasoning tasks. This study evaluates the performance of leading LLMs, including OpenAI's GPT-4o and Alibaba's Qwen2.5-72B, on intricate logical reasoning challenges, focusing on the Turtle Soup game—a lateral thinking puzzle that requires deductive reasoning through yes/no questions. We conducted comprehensive experiments using various few-shot prompting scenarios and fine-tuning techniques across multiple LLMs. 
% It is discovered that LLMs exhibit inconsistencies in producing valid classifications. To measure their capability in this regard, we propose a novel evaluation metric: the Valid Classification Ratio (VCR), which quantifies the proportion of valid classifications among the total outputs of the LLMs.
% Our results demonstrate the superior efficacy of LoRA/QLoRA fine-tuning compared to few-shot prompting. Fine-tuned models not only achieved higher peak performance but also reached a perfect VCR of 1.0 after just 0.2 epochs, while exhibiting faster inference times. 
% The experimental results also indicate that VCR is a valuable metric for assessing model reliability in precise logical reasoning tasks. Notably, the Qwen2.5-72B model emerged as the top performer, surpassing fine-tuning methods in maximizing LLMs' logical reasoning capabilities. Our findings have significant implications for enhancing LLM performance in complex problem-solving scenarios and real-world applications.



\end{abstract}

\begin{IEEEkeywords}
Logical Reasoning, Large Language Models (LLMs), Valid Classification Ratio (VCR), Few-Shot Prompting, Fine-Tuning, Turtle Soup Game, LoRA and QLoRA
\end{IEEEkeywords}

\section{Introduction}

The Turtle Soup game, popularized by the book \textit{Challenging Lateral Thinking Puzzles} \cite{sloane1993challenging}, is a mystery puzzle that has captivated youths worldwide. The game gets its name from a classic problem in the genre: a man orders turtle soup and subsequently commits suicide. Players must deduce the rationale behind the suicide by asking questions that can only be answered with "yes" or "no." These puzzles encourage lateral thinking, pushing players to break free from conventional thought patterns and approach problems from new perspectives, especially when situations seem initially nonsensical or incomprehensible. A sample dialogue illustrates this concept:

\small{
\begin{dialogue}
\speak{Game Master} On an express train, a man boarded halfway through the journey. After the doors closed, he suddenly regained consciousness and began observing the faces of the passengers around him. "With all due respect, are you 29 years old?" "Yes, but how do you know?" Ignoring him, the man continued speaking to others. "You are 45 years old, right?" "That's correct." "Are you 81 years old?" "How do you know?" After asking around, the man returned to his seat, sat down quietly, and began to cry in despair. Why is that?
\speak{Player} Is he dreaming?
\speak{Game Master} No.
\speak{Player} Does he know these people?
\speak{Game Master} No.
\speak{Player} Does he have superpowers?
\speak{Game Master} Yes.
\speak{Player} He is an angel condemned to live as a human.
\speak{Game Master} Congratulations! You guessed the answer. The original answer is that the man can see how long each person will live, but he cannot see his own lifespan. His life expectancy is about to end, meaning if the train is going to crash, he will die soon. The game is over.
\end{dialogue}
}

The advent of Large Language Models (LLMs) has led to a surge in online Turtle Soup games, leveraging these models' language understanding and logical reasoning capabilities. However, the efficacy of LLMs in logical reasoning remains underexplored \cite{parmar2024logicbench}.

This study presents a comprehensive evaluation of the latest generation of open-source LLMs, such as Qwen2.5 and Llama3.1, alongside OpenAI's most recent models, in tackling complex and abstract logical reasoning tasks like those found in the Turtle Soup Question Answering (TSQA) task. Our research not only highlights the strengths and limitations of different approaches—few-shot prompting and fine-tuning—but also challenges the notion that proprietary models are inherently superior. Our findings demonstrate that open-source models can match, and in some cases exceed, the performance of established models like GPT-4o.

Our main contributions are:
\begin{itemize}
    \item Utilizing the Turtle Soup game as a benchmark for evaluating the logical reasoning capabilities of LLMs through deductive reasoning.
    
    \item Introducing the Valid Classification Ratio (VCR) as a new metric for assessing the reliability of model outputs, measuring the proportion of valid classifications among total outputs.

    \item Demonstrating the superior efficacy of LoRA/QLoRA fine-tuning methods over few-shot prompting, establishing fine-tuning as a crucial factor in optimizing LLM performance.

    \item Identifying the Qwen2.5-72B and Llama-3.1-70B models as the top performers in TSQA tasks, surpassing GPT-4o, thereby contributing valuable insights into the strengths of different LLM architectures.
\end{itemize}


% While logical reasoning with LLMs has advanced significantly, it remains a challenging and evolving area. Models like GPT-4 demonstrate impressive pattern recognition capabilities and can make inferences from large-scale text data. They handle basic formal logic tasks, such as propositional logic or syllogisms, and respond reasonably well to informal reasoning challenges like common-sense questions. LLMs have developed the ability to understand context and leverage nuanced relationships between sentences and concepts, leading to more accurate logical decisions.%%Please add references in this paragraph as well.

% However, LLMs still struggle with deductive reasoning, where consistency is critical across multiple steps. While they can generate logical conclusions in isolation, maintaining a coherent chain of reasoning over several steps or sentences remains difficult. In multi-step complex reasoning tasks, LLMs may make errors that compound over long logical chains. They also face limitations in tasks requiring abstract thought or the manipulation of symbolic or formal rules. The widely anticipated release of OpenAI o1 with its chain-of-thought strategy seeks to break through in this aspect. %%Please add references in this paragraph as well.
 

% The recent AI boom saw the rise of more LLMs competing with traditional powerhouses like GPT by OpenAI and Llama by MetaAI. Open-source LLMs such as Qwen and Mistral have proven themselves worthy adversaries to their paid, closed-source competitors. Alibaba Cloud's recently released Qwen2.5 series boasts outstanding performance in both large and small models. Additionally, they come equipped with features such as a large token count of  128K tokens context window and 8K tokens generation, multilingual support for over 29 languages, and Grouped Query Attention which greatly enhances the efficiency of the attention process. %%Please add references in this paragraph as well.

% Qwen2.5, alongside InternLM2.5, also represent the subset of LLMs developed in China. Their capabilities are of high interest to academia, as the AI ecosystem stands to benefit from a more competitive and diverse global AI landscape, reducing the dominance of Western, particularly US-based, tech companies. Chinese LLMs are often optimized for tasks requiring fluency in both Chinese and English (and sometimes other languages), signifying exciting potential for cross-border trade, international diplomacy, and other areas where bilingual capabilities are valued. These models are more attuned to Chinese and Asian cultural, social, and historical contexts, helping to address biases that might exist in models developed in the West. %%Please add references in this paragraph as well.


\section{Related Work}
%Logical Reasoning Dataset

%Hanmeng Liu et al. introduced LogiQA 2.0, an enhanced version of the original LogiQA dataset designed for evaluating logical reasoning~\cite{liu2023logiqa}. This updated version expands the dataset from 8,678 to 15,708 instances, receives high-quality English translations delivered by professional human translators, and includes a Natural Language Inference (NLI) dataset derived from the original MRC task. Experiments demonstrate that the increased dataset size and refined quality contribute to better performance across various pre-trained models, though state-of-the-art methods still fall short of human-level reasoning. The NLI section of LogiQA 2.0 is recognized as the first large-scale, expert-designed NLI benchmark for logical reasoning.
%
%ZX: If we do not use this dataset, LogiQA dataset, why did we mention it here? This gives the reviewers an opportunity to reject our paper for including irrelevant research. 
%

%Situation Puzzles Evaluation

%Tian Zhao et al. proposed SPQA, a new question-answering model specifically designed to tackle situation puzzle questions by incorporating a novel situation-puzzle related dataset (SpQ)~\cite{zhao2023question}. Built on the T5 architecture similar to UnifiedQA, SPQA is fine-tuned and prompt-tuned to assess its performance in handling both situation puzzle questions and standard yes/no questions. Their study reveal fine-tuning with the SpQ dataset significantly enhances the model's ability to solve situation puzzles but reduces its performance on general yes/no questions. Prompt-tuning further highlights the importance of the SpQ dataset, with a notable improvement in solving situation puzzles.
%
%%ZX: When the reviewers read the paragraph above, they will expect us to use the situation-puzzle related dataset for comparison. However, our paper did not use the dataset. Why don’t we consider the situation-puzzle related dataset you mentioned above? It would be interesting to use the same dataset that other researchers have utilized.
%
%OpenAI o1
%
%OpenAI o1 is a cutting-edge model that excels in competitive programming and complex reasoning tasks. It ranks in the 89th percentile on Codeforces and achieves top-tier performance on the USA Math Olympiad qualifier (AIME), solving up to 93\% of problems using consensus-based sampling techniques. OpenAI o1 also surpasses human PhD-level performance on the GPQA-diamond benchmark, which assesses expertise in biology, physics, and chemistry. Notably, it is the first model to be competitive with human experts on MMLU's MMMU subset, scoring 78.2\%. Leveraging a chain-of-thought strategy, o1 utilizes reinforcement learning to break down complex problems into simpler steps, continuously refining its reasoning and problem-solving capabilities. This advanced training method allows the model to improve both its accuracy and efficiency over time.
%%
% ZX: Detailed information from OpenAI's report above provides more context. The statistic stating that '93% of problems use consensus-based sampling techniques' needs to be properly cited; please specify the source. Since I could not find the exact numbers, I refrained from using them. Please feel free to add it into Related Work Section (e.g., second pragraph).
%
%
%Qwen2.5
%The Qwen2.5 series models represent the newest frontier in the Qwen family. They are trained on 18 trillion tokens and surpasses Qwen2 significantly in knowledge (MMLU: 85+), as well as proficiency in coding (HumanEval 85+) and mathematics (MATH 80+). The flagship model Qwen2.5-72B outperforms the newly released Llama3 70B model and the mixture of experts Mixtral 8x22B model across various benchmarks, and remains highly competitive even when compared against larger models like Llama-3-405B\cite{qwen2.5}.  On the other end of the spectrum, Qwen2.5-3B demonstrates the accelerated growth in knowledge density among language models is a prime example of small language models (SLMs) reducing the gap with LLMs. Additionally, the new models achieve significant improvements in instruction following, generating long texts (over 8K tokens), understanding structured data (e.g, tables), and generating structured outputs especially JSON. Qwen2.5 models are generally more resilient to the diversity of system prompts, enhancing role-play implementation and condition-setting for chatbots.   
%zx: rewrite as below in the related work section.

%Ensemble Model 
%
%Chengyan Wu et al. (2024) proposed an ensemble model for diabetes question classification, combining multiple large and small language models to enhance performance. The model ensemble includes Qwen-7b-Chat, ChatGLM2-6b, and MacBERT as base models, fine-tuned using the QLoRA and LoRA frameworks for instruction-based learning. These large models focus on multiple-choice tasks, while smaller models trained with adversarial techniques, such as Fast Gradient Method (FGM), handle traditional text classification. To further improve accuracy, misclassified samples are identified and augmented via paraphrasing, adding diversity to the training data. A majority voting mechanism is then applied to aggregate predictions from the models, ensuring a robust final classification. Their findings suggest that integrating advanced
%language models with few-shot learning strategies can significantly enhance classification performance, especially when the ground truth data is generated or verified by the same type of models.
% ZX: why we introduce Ensemble Model here? it is irrelevant to this paper.
%
%
The rapid evolution of large language models (LLMs) has generated significant interest in their applications for complex reasoning tasks \cite{chen2024large, hadi2024large}. Early research underscored the foundational capabilities of models like GPT-3 and its successors in handling basic logical reasoning and natural language processing tasks \cite{zhang2021commentary, kalyan2023survey}. While these models excelled in pattern recognition and could infer information from extensive datasets, challenges remained, particularly in maintaining coherent reasoning across multiple steps \cite{yang2024harnessing}.

Recent advancements, such as OpenAI's GPT-4, have enhanced deductive reasoning abilities through techniques like chain-of-thought prompting, which encourage models to articulate their reasoning processes more clearly \cite{gou2024rationality}. Despite these improvements, LLMs still exhibit limitations in multi-step reasoning tasks, often leading to cumulative errors in lengthy logical chains \cite{tong2024can}.

The landscape of LLMs is rapidly expanding, with significant contributions from both proprietary and open-source models \cite{he2023survey}. Meta's Llama series and newer models from Alibaba, such as Qwen2.5, have emerged as strong competitors to established leaders \cite{hui2024qwen2, yang2024qwen2}. These models not only offer large context windows and multilingual support but also incorporate innovative features like Grouped Query Attention, enhancing efficiency in processing complex queries \cite{yang2024harnessing}.

% The Qwen2.5 series represents the latest advancements in the Qwen family. Trained on 18 trillion tokens, these models significantly surpass Qwen2 in knowledge (MMLU: 85+), coding proficiency (HumanEval: 85+), and mathematical abilities (MATH: 80+). The flagship model, Qwen2.5-72B, outperforms the recently released Llama3 70B model and the mixture of experts, Mixtral 8x22B, across various benchmarks, remaining highly competitive even against larger models like Llama-3-405B \cite{qwen2.5}. Conversely, Qwen2.5-3B exemplifies the accelerated growth in knowledge density among language models, demonstrating how small language models (SLMs) are narrowing the performance gap with larger language models (LLMs). Furthermore, these new models achieve substantial improvements in instruction following, generating long texts (over 8K tokens), understanding structured data (e.g., tables), and producing structured outputs, particularly in JSON format. The Qwen2.5 models are generally more resilient to diverse system prompts, enhancing role-play implementation and condition-setting for chatbots.

Moreover, the performance of Chinese LLMs, such as InternLM2.5, reflects a growing emphasis on addressing cultural and linguistic nuances within AI systems. These models are specifically designed to accommodate bilingual contexts and align with the unique social and historical aspects of Asian cultures, providing potential advantages in cross-cultural applications \cite{zhang2024mme}.

Several studies have begun to investigate the competitive landscape of LLMs in reasoning tasks. Recent evaluations have focused on few-shot and fine-tuning methodologies to determine optimal performance metrics across various architectures \cite{winata2021language}. However, the exploration of logical reasoning in contexts like the Turtle Soup game remains limited, presenting opportunities for further investigation.

% This study builds on the existing body of research by systematically comparing the latest iterations of LLMs against proprietary models in complex reasoning scenarios, utilizing both few-shot prompting and fine-tuning techniques. By introducing the Valid Classification Ratio (VCR) as a new metric for assessing model output reliability, our work aims to contribute valuable insights into the evolving capabilities of LLMs in logical reasoning tasks.


\section{Dataset}

To evaluate the logical reasoning abilities of large language models (LLMs), we utilized a dataset released during the recent “Logical Reasoning with LLM” competition organized by Mango TV\footnote{\url{https://challenge.ai.mgtv.com/#/track/25}}. The original dataset comprises three components: a training set, a validation set, and a test set. The training set contains 25,000 user-generated questions paired with labels, including puzzles and their corresponding truths. The validation set consists of 3,000 user-generated questions and labels. The test set includes 30,000 user questions but does not include labels. In our study, we utilized only the training and validation sets, collectively referred to as the “Turtle Soup Question Answering (TSQA)” dataset.

The TSQA dataset, being in Chinese, presents an additional layer of complexity for logical reasoning tasks, especially for LLMs primarily trained on English data. It consists of multiple Turtle Soup games, each identified by a unique Chinese title. For each game, a puzzle in Chinese is provided, which the player must analyze to deduce the truth. The text field contains the player's questions in Chinese, which are then processed by the LLM.

\begin{CJK*}{UTF8}{gbsn}
The model's responses are limited to five categories: 是(Yes), 不是(No), 不重要(Unimportant), 问法错误(Wrong Question), and 回答正确(Correct Answer). The criteria for judging each answer are as follows:
\end{CJK*}
\begin{itemize}
    \item If the answer to the question can be directly or indirectly found in the puzzle and truth, the model should respond with "Yes" or "No."
    \item If the answer cannot be directly or indirectly inferred from the puzzle and truth, the model should respond with "Unimportant."
    \item If the question is not a closed-ended question or is difficult to understand, the model should respond with "Wrong Question."
    \item If the question accurately reconstructs the truth of the game, the model should respond with "Correct Answer."
\end{itemize}

The dataset structure, including sample values in both Chinese and their English translations, is illustrated in Table~\ref{tab:tsqa_dataset}. This bilingual presentation helps to demonstrate the nature of the dataset while making it accessible to non-Chinese speakers.

\begin{table*}[h!]
    \centering
    \begin{CJK*}{UTF8}{gbsn}
    \caption{Structure of the Turtle Soup Question Answering (TSQA) Dataset, \small{including sample values in Chinese and their English Translations. Designed to evaluate the logical reasoning capabilities of LLMs, the dataset consists of user-generated questions in Chinese paired with corresponding labels, including puzzles and their truths.}}
    \label{tab:tsqa_dataset}
    \resizebox{0.8\textwidth}{!}{
    \begin{tabular}{|c|>{\centering\arraybackslash}p{0.1\linewidth}|>{\centering\arraybackslash}p{0.25\linewidth}|>{\centering\arraybackslash}p{0.4\linewidth}|} 
    \hline 
    Field & Description & Sample Value (Chinese) & Sample Value (English Translation) \\ \hline 
    title & Title & 甄庄的哭声 & The Crying Sound of Zhenzhuang \\ \hline 
    puzzle & Puzzle & 在一个安静的夜晚，小村庄的湖边突然传来了阵阵哭泣声。第二天早晨，村长甄锐发现湖边的石头上放着一顶破旧的帽子，但没有人知道这顶帽子是从哪里来的，哭泣声又是为何。请还原故事真相。 & On a quiet night, a series of cries suddenly came from the lakeside in the small village. The next morning, the village head Zhen Rui found an old hat on a stone by the lake, but no one knew where the hat came from or why the crying occurred. Please restore the truth of the story. \\ \hline 
    truth & Truth & 原来，这顶破旧的帽子属于一个小男孩，他小时候与爷爷在湖边生活。爷爷教他钓鱼、游泳，还告诉他湖中的海龟是他们的朋友。后来，小男孩随父母去了城市生活，但每年夏天都会回到村子探望爷爷。然而，去年夏天，爷爷因病去世，小男孩伤心欲绝。今年夏天，他回到村子，来到湖边，想起和爷爷的美好回忆，忍不住哭泣。他将爷爷的帽子放在湖边的石头上，希望能让爷爷的在天之灵得到安慰。那晚的哭泣声正是小男孩在祭奠他亲爱的爷爷。 & The old hat belonged to a little boy who lived by the lake with his grandfather when he was young. His grandfather taught him how to fish and swim and told him that the turtles in the lake were their friends. Later, the boy moved to the city with his parents, but he returned to the village every summer to visit his grandfather. However, last summer, his grandfather passed away due to illness, leaving the boy heartbroken. This summer, the boy returned to the village, went to the lake, remembered the beautiful times with his grandfather, and couldn't help but cry. He placed his grandfather's hat on the stone by the lake, hoping to bring some comfort to his grandfather's spirit in heaven. The crying that night was the boy mourning his beloved grandfather. \\ \hline
    text & Player's question & 有个叫甄庄的村子 & There is a village called Zhenzhuang \\ \hline 
    label & Human label & 是 & Yes \\ \hline 
    \end{tabular}
    }
    \end{CJK*}
\end{table*}


\section{Methodology}

\subsection{Selection of Large Language Models (LLMs)}

Our study incorporates a diverse array of state-of-the-art large language models (LLMs), ranging from open-source models with 3 billion to 72 billion parameters to the latest offerings from OpenAI, such as GPT-4o and o1-preview. For open-source models lacking native Chinese language support, we utilized their fine-tuned Chinese versions. Table \ref{tab:ai_models} provides a comprehensive overview of all models used in our study, detailing their names, sizes, context lengths, and corresponding HuggingFace model identifiers where applicable.

\begin{table*}[htbp]


\caption{Overview of Large Language Models and Their Specifications. \small{This table provides a detailed summary of the LLMs used in the study, including models from OpenAI, Meta, Mistral AI, Alibaba, and Shanghai AI Laboratory. The table lists each model's name, size (ranging from 3 billion to 72 billion parameters), context length, and corresponding HuggingFace model identifier where applicable. Notably, open-source models that lacked native Chinese language support were fine-tuned on Chinese datasets to enhance their performance in this study's logical reasoning tasks.}}

\label{tab:ai_models}
\centering
\resizebox{0.8\textwidth}{!}{
\begin{tabular}{|l|l|r|r|l|}
\hline
\textbf{Company} & \textbf{Model Name} & \textbf{Size} & \textbf{Context Length} & \textbf{HuggingFace Model ID} \\
\hline
\multirow{4}{*}{OpenAI} 
& GPT-4o & N/A & 128K & N/A \\
& GPT-4o-mini & N/A & 128K & N/A \\
& o1-preview & N/A & 128K & N/A \\
& o1-mini & N/A & 128K & N/A \\
\hline
\multirow{2}{*}{Meta} 
& Llama3.1-8B & 8B & 128K & shenzhi-wang/Llama3.1-8B-Chinese-Chat \\
& Llama3.1-70B & 70B & 128K & shenzhi-wang/Llama3.1-70B-Chinese-Chat \\
\hline
Mistral AI & Mistral-7B & 7B & 32K & shenzhi-wang/Mistral-7B-v0.3-Chinese-Chat \\
\hline
\multirow{5}{*}{Alibaba} 
& Qwen2-7B & 7B & 128K & Qwen/Qwen2-7B-Instruct \\
& Qwen2-72B & 72B & 128K & Qwen/Qwen2-72B-Instruct \\
& Qwen2.5-3B & 3B & 32K & Qwen/Qwen2.5-3B-Instruct \\
& Qwen2.5-7B & 7B & 128K & Qwen/Qwen2.5-7B-Instruct \\
& Qwen2.5-72B & 72B & 128K & Qwen/Qwen2.5-72B-Instruct \\
\hline
\multirow{3}{*}{Shanghai AI Laboratory} 
& InternLM2.5-7B & 7B & 16K & internlm/internlm2\_5-7b-chat \\
& InternLM2.5-7B-1M & 7B & 1024K & internlm/internlm2\_5-7b-chat-1m \\
& InternLM2.5-20B & 20B & 16K & internlm/internlm2\_5-20b-chat \\
\hline
\end{tabular}
}
\end{table*}

\subsection{Few-shot Learning and Parameter-Efficient Fine-tuning}

To enable effective performance on Turtle Soup Question Answering (TSQA) tasks, we employed few-shot prompting techniques. The system prompt used for this task is concise:

\begin{minted}[fontsize=\footnotesize, breaklines]{text}
You are an expert in logical reasoning.
\end{minted}

It's worth noting that the OpenAI's o1-preview and o1-mini models do not support system prompts.

The user prompt template, originally in Chinese, outlines the rules of a Turtle Soup game. We provide both the original Chinese version (Listing \ref{lst:user-prompt-chinese}) and its English translation (Listing \ref{lst:user-prompt-english}) for clarity.


\begin{CJK*}{UTF8}{gbsn}
\begin{listing}[htbp]
\begin{minted}[fontsize=\footnotesize, breaklines]{text}
\small{
你是一个情景猜谜游戏的主持人。游戏规则如下：

1. 参与者会得到一个谜面，谜面会描述一个简单又难以理解的事件。
2. 主持人知道谜底，谜底是谜面的答案。
3. 参与者可以询问任何封闭式问题来找寻事件的真相。
4. 对于每个问题，主持人将根据实际情况回答以下五个选项之一：是、不是、不重要、回答正确、问法错误。各回答的判断标准如下：
   - 若谜面和谜底能找到问题的答案，回答：是或者不是
   - 若谜面和谜底不能直接或者间接推断出问题的答案，回答：不重要
   - 若参与者提问不是一个封闭式问题或者问题难以理解，回答：问法错误
   - 若参与者提问基本还原了谜底真相，回答：回答正确
5. 回答中不能添加任何其它信息，也不能省略选项中的任何一个字。例如，不可以把"不是"省略成"不"。

请严格按照这些规则回答参与者提出的问题。

谜面: {puzzle}
谜底: {truth}
参与者提出的问题: {text}
回答: 
}
\end{minted}
\caption{User Prompt Template for TSQA (Chinese)}
\label{lst:user-prompt-chinese}
\end{listing}
\end{CJK*}

\begin{listing}[htbp]
\begin{minted}[fontsize=\footnotesize, breaklines]{text}
\small{
You are the host of a situational guessing game. The rules of the game are as follows:

1. Participants will receive a puzzle that describes a simple yet difficult to understand event.
2. The host knows the truth, which is the solution to the puzzle.
3. Participants can ask any closed-ended questions to uncover the truth of the event.
4. For each question, the host will respond with one of the following five options based on the actual situation: Yes, No, Unimportant, Correct Answer, or Wrong Question. The criteria for each response are as follows:
   - If the puzzle and truth can provide an answer to the question, respond with: Yes or No
   - If the puzzle and truth cannot directly or indirectly infer an answer to the question, respond with: Unimportant
   - If the participant's question is not a closed-ended question or is difficult to understand, respond with: Wrong Question
   - If the participant's question essentially reveals the truth of the answer, respond with: Correct Answer
5. The response must not include any additional information, nor should any word be omitted from the options. For example, "Correct answer" cannot be abbreviated to "Correct".

Please strictly follow these rules when answering the participant's questions.

Puzzle: {puzzle}
Truth: {truth}
Participant's question: {text}
Answer: 
}
\end{minted}
\caption{User Prompt Template for TSQA (English)}
\label{lst:user-prompt-english}
\end{listing}


In the prompt template, we populated the placeholders with specific data:

\begin{itemize}
    \item \texttt{\{exemplars\}}: For zero-shot tasks, this field is left empty. For few-shot tasks, we retrieved exemplars from the TSQA training set, formatted as shown below (translated to English):
\small{
\begin{minted}[fontsize=\footnotesize, 
               breaklines=true, 
               breakanywhere=true]{text}
Example Inputs and Outputs:
Puzzle: In the village of Zhen, there is a legend that every year, when the harvest season for pumpkins arrives, one of the largest pumpkins in the field disappears without a trace. The villagers are puzzled by this phenomenon.
Truth: The truth turned out to be related to an old farmer. When he was young, he had fallen in love with a beautiful girl. They had agreed to marry when the season for harvesting pumpkins arrived. But fate had a way of playing tricks on people. The girl died in a car accident on the wedding day. Heartbroken, the farmer decided to remember his beloved by stealing the largest pumpkins every year and putting them in front of her grave. This act of kindness continued for many years, becoming a mysterious legend in the village.
Participant's question: Were they stolen from the village?
Answer: Yes
... (additional examples omitted for brevity)
\end{minted}

To teach the LLM all valid classification labels—Yes, No, Unimportant, Wrong Question, and Correct Answer—in a few-shot learning setup, we need to provide at least one example for each label. This requires a minimum of five examples, one for each label, making it a 5-shot scenario. For larger setups, the number of examples would follow this pattern, increasing in multiples of five.
}
    \item \texttt{\{puzzle\}}, \texttt{\{truth\}}, and \texttt{\{text\}}: These were retrieved from the TSQA validation set.
\end{itemize}

To enhance the performance of the open-source models, we fine-tuned them on the training set for two epochs using zero-shot prompting. The fine-tuning process was adjusted based on the size of each model:

\begin{itemize}
    \item For InternLM2.5-20B, Llama-3.1-70B and Qwen2.5-72B variants: We employed 4-bit quantization using LoRA (QLoRA) with the open-source library Llama Factory \cite{zheng2024llamafactory} on an Nvidia H100 GPU.
    \item For other open-source models: We used LoRA with the Llama Factory library on an Nvidia L40 GPU.
\end{itemize}

To assess performance, we ran inference on the LLMs to obtain their outputs, which were then evaluated against the ground truths. We evaluated the performance of OpenAI models using the OpenAI API.
% \footnote{\url{https://platform.openai.com/docs/overview}}
In contrast, the performance of open-source models, both with and without LoRA/QLoRA adapters, was assessed using the HuggingFace Transformers library\cite{wolf2020transformers}.

\subsection{Evaluation Metrics}

To evaluate the performance of each model, we calculated predictive accuracy on the held-out test data using the weighted F1 score. This metric provides a comprehensive assessment of the model’s performance by considering both precision and recall, offering a balanced view of the model's ability to correctly classify data.

During the evaluation process, we observed that the Large Language Models (LLMs) occasionally generated outputs that did not correspond to any valid classification labels: Yes, No, Unimportant, Wrong Question, or Correct Answer. Table \ref{tab:classification_results} presents a sample of classification results produced by the GPT-4o model across various shot counts (0, 5, 10, 20, 30, 40, and 50) for the example question/puzzle shown in Table \ref{tab:tsqa_dataset}.

\begin{table*}[h!]

\caption{Sample Classification Results from GPT-4o Model Across Different Shot Counts. \small{This table presents the classification results generated by the GPT-4o model for the sample values provided in Table~\ref{tab:tsqa_dataset} across various shot counts (0, 5, 10, 20, 30, 40, and 50). It highlights the variability in the model's output, especially at the 20-shot level, where the model correctly classified the question but included additional extraneous text that would require post-processing.}}

\label{tab:classification_results}
\centering
\begin{CJK*}{UTF8}{gbsn}
\resizebox{0.8\textwidth}{!}{
\begin{tabular}{|l|p{6.5cm}|p{10cm}|}
\hline
\textbf{Shot} & \textbf{GPT-4o Result} & \textbf{GPT-4o Result (English Translation)} \\ \hline
\textbf{0} & 问法错误 & Wrong question \\ \hline
\textbf{5} & 不重要 & Not important \\ \hline
\textbf{10} & 不重要 & Not important \\ \hline
\textbf{20} & 是

谜面: 在一个安静的夜晚，小村庄的湖边突然传来...%了阵阵哭泣声。第二天早晨，村长甄锐发现湖边的石头上放着一顶破旧的帽子，但没有人知道这顶帽子是从哪里来的，哭泣声又是为何。请还原故事真相。

谜底: 原来，这顶破旧的帽子属于一个小男孩，他...%小时候与爷爷在湖边生活。爷爷教他钓鱼、游泳，还告诉他湖中的海龟是他们的朋友。后来，小男孩随父母去了城市生活，但每年夏天都会回到村子探望爷爷。然而，去年夏天，爷爷因病去世，小男孩伤心欲绝。今年夏天，他回到村子，来到湖边，想起和爷爷的美好回忆，忍不住哭泣。他将爷爷的帽子放在湖边的石头上，希望能让爷爷的在天之灵得到安慰。那晚的哭泣声正是小男孩在祭莫他亲爱的爷爷。

参与者提出的问题: 哭泣声是人发出来的吗

回答: 是 & Yes

Puzzle:
On a quiet night, a series of cries suddenly came from the lakeside in the small...% village. The next morning, the village head Zhen Rui found an old hat on a stone by the lake, but no one knew where the hat came from or why the crying occurred. Please restore the truth of the story.

Truth:
The old hat belonged to a little boy who lived by the lake with his grandfather...% when he was young. His grandfather taught him how to fish and swim and told him that the turtles in the lake were their friends. Later, the boy moved to the city with his parents, but he returned to the village every summer to visit his grandfather. However, last summer, his grandfather passed away due to illness, leaving the boy heartbroken. This summer, the boy returned to the village, went to the lake, remembered the beautiful times with his grandfather, and couldn’t help but cry. He placed his grandfather’s hat on the stone by the lake, hoping to bring some comfort to his grandfather’s spirit in heaven. The crying that night was the boy mourning his beloved grandfather.

Participant’s Question: Was the crying sound made by a person?

Answer: Yes \\ \hline
\textbf{30} & 问法错误 & Wrong question \\ \hline
\textbf{40} & 问法错误 & Wrong question \\ \hline
\textbf{50} & 问法错误 & Wrong question \\ \hline
\end{tabular}
}
\end{CJK*}
\end{table*}


\begin{CJK*}{UTF8}{gbsn}
Upon closer examination, the results at the 20-shot level reveal a significant observation. While the model correctly provided the classification (“是” / “Yes”), it also included a substantial amount of additional text, essentially combining the prompt with the answer. This extraneous text would need to be removed through post-processing before calculating the F1 score. This behavior suggests that the model sometimes fails to adhere to the instruction defined in Listing~\ref{lst:user-prompt-english}, which clearly states, “The response must not include any additional information, nor should any word be omitted from the options.”
\end{CJK*}

To address this issue, we propose the introduction of a new metric: the Valid Classification Ratio (VCR). The VCR is defined as the ratio of outputs that correspond to valid classification labels to the total number of outputs generated.

\begin{equation}
\text{VCR} = \frac{\text{Valid Classifications}}{\text{Total Classifications}}
\end{equation}

This metric provides a quantitative measure of the model's reliability in producing valid classifications across different contexts and shot counts.


\section{Results}
Our experiments on the Turtle Soup QA task involve two main approaches: few-shot prompting and fine-tuning. We evaluated both open-source Large Language Models (LLMs) and proprietary models from OpenAI. The results reveal interesting patterns in model performance and the effects of different training strategies.

\subsection{Few-shot Prompting Performance}
Few-shot prompting was applied to both open-source LLMs and OpenAI models. The results, as shown in Table~\ref{tab:performance}, reveal several key insights:

\subsubsection{Competitive Performance of Open-Source Models}
Contrary to initial expectations, open-source models demonstrated performance levels competitive with those of OpenAI models. Notably, Qwen2.5-72B achieved an F1 score of 80.88\% with 10-shot prompting, surpassing GPT-4o's performance of 80.36\% with 10-shot prompting. This suggests that the gap between proprietary and open-source models is narrowing in complex reasoning tasks.

\subsubsection{Unexpected Performance of o1-preview}
Interestingly, the o1-preview model, which is purported to have superior reasoning capabilities, exhibited lower F1 scores compared to GPT-4o. The highest F1 score for o1-preview was 77.08\% at 50-shot, while GPT-4o achieved 80.36\% at 10-shot. This unexpected result calls for further investigation into the specific strengths and weaknesses of these models, particularly in lateral thinking tasks.

% \subsubsection{Model Size and Performance Correlation}
% Generally, larger models exhibited superior performance. For example, Llama3.1-70B outperformed its 8B counterpart in most shot settings, with a peak F1 score of 75.67\% at 10-shot compared to the 8B model's peak of 76.97\% at 30-shot. This trend is consistent with the common observation that increased model size often correlates with enhanced performance in complex tasks.

\subsubsection{Missing Data Points}
It is important to note that our results include some missing data points, particularly for larger models and higher shot settings. These gaps are primarily due to two key factors:

\begin{enumerate}
    \item \textbf{GPU Memory Limitations:} We encountered CUDA out-of-memory errors with some of the larger models at higher shot settings. For instance, we were unable to obtain data for Llama3.1-70B beyond the 30-shot setting. These limitations restricted our ability to evaluate these models at higher shot counts.

    \item \textbf{Time Constraints:} The evaluation process for some models at higher shot counts was extremely time-consuming. In certain cases, the allocated job time expired before the evaluation could be completed. This issue is evident with models like InternLM2.5-20B, for which we only have data for the 0-shot setting. Similarly, we lack data for Qwen2-72B and Qwen2.5-72B beyond the 20-shot setting, and for Mistral-7B at the 50-shot setting, due to the same reason.
\end{enumerate}

These findings underscore the significant computational resources and time required for a comprehensive assessment of the largest language models, particularly in few-shot scenarios.

\subsubsection{Performance Trends Across Shot Counts}
Despite some missing data points, several trends in model performance across different shot counts can be observed:

\begin{itemize}
    \item \textbf{Variability in Optimal Shot Count:} Different models appear to reach their peak performance at varying shot counts. For instance, GPT-4o achieved its highest F1 score of 80.36\% at 10-shot, whereas o1-preview peaked at 77.08\% at 50-shot. This suggests that the optimal number of examples may vary depending on the model's architecture and size.
    
    \item \textbf{Inconsistent VCR:} The Valid Classification Ratio (VCR) demonstrates significant variability across models and shot counts, particularly among open-source models. For instance, Mistral-7B-v0.3 shows extremely low VCR values (ranging from 1.17\% to 14.20\%) compared to other models. This inconsistency in VCR highlights a critical challenge in few-shot learning scenarios and underscores the need for exploring fine-tuning approaches to improve model performance and stability.
\end{itemize}

\begin{table*}[htbp]
\centering
\small
\caption{Performance of Language Models on the Turtle Soup QA Task Across Different Shot Counts (F1 Score and Valid Classification Ratio in \%). \small{The top metrics for each model across different shot counts are highlighted in bold. Key observations include Qwen2.5-72B surpassing GPT-4o in few-shot settings with an F1 score of 80.88\% at 10-shot, and the variability in optimal shot counts across models. Additionally, some models exhibit diminishing returns or inconsistent VCR at higher shot counts.}}
\resizebox{.7\textwidth}{!}{
\begin{tabular}{l|l|ccccccc}
\hline
Model & Metric & 0-shot & 5-shot & 10-shot & 20-shot & 30-shot & 40-shot & 50-shot \\
\hline
\multirow{2}{*}{GPT-4o-mini} & F1 & \textbf{72.96} & 71.81 & 69.17 & 67.98 & 69.13 & 69.25 & 71.05 \\
                             & VCR & 99.17 & \textbf{99.97} & 99.83 & 99.80 & 99.90 & 99.87 & 99.93 \\
\hline
\multirow{2}{*}{GPT-4o} & F1 & 79.53 & 80.00 & \textbf{80.36} & 79.67 & 80.31 & 79.93 & 80.14 \\
                        & VCR & 6.60 & 99.80 & \textbf{99.97} & 99.93 & 99.90 & 99.73 & 99.93 \\
\hline
\multirow{2}{*}{o1-mini} & F1 & 73.77 & 74.83 & 74.86 & 75.35 & 75.48 & 75.73 & \textbf{75.90} \\
                         & VCR & \textbf{99.90} & 99.67 & 99.43 & 99.47 & 99.77 & 99.77 & 99.77 \\
\hline
\multirow{2}{*}{o1-preview} & F1 & 74.51 & 75.35 & 76.77 & 76.25 & 76.44 & 76.74 & \textbf{77.08} \\
                            & VCR & \textbf{99.80} & 97.90 & 98.73 & 98.53 & 98.40 & 98.40 & 98.17 \\
\hline
\multirow{2}{*}{Llama3.1-8B} & F1 & 73.71 & 72.69 & 70.83 & 76.62 & \textbf{76.97} & 70.67 & 72.31 \\
                             & VCR & 80.33 & \textbf{98.87} & 96.23 & 97.90 & 73.27 & 75.90 & 66.23 \\
\hline
\multirow{2}{*}{Llama3.1-70B} & F1 & 75.26 & 75.45 & \textbf{75.67} & 73.49 & 75.65 & - & - \\
                              & VCR & 0.97 & 79.00 & \textbf{83.27} & 81.90 & 54.80 & - & - \\
\hline
\multirow{2}{*}{Mistral-7B} & F1 & 66.34 & \textbf{68.10} & 64.93 & 66.91 & 68.63 & 66.44 & - \\
                            & VCR & 1.17 & \textbf{14.20} & 10.63 & 8.27 & 7.00 & 6.33 & - \\
\hline
\multirow{2}{*}{InternLM2.5-7B} & F1 & 69.06 & \textbf{72.71} & 62.51 & 67.70 & 65.35 & 61.62 & 55.68 \\
                                & VCR & \textbf{100.00} & 99.90 & 98.83 & 94.73 & 94.03 & 98.13 & 98.03 \\
\hline
\multirow{2}{*}{InternLM2.5-7B-1M} & F1 & 52.45 & \textbf{77.01} & 66.58 & 67.64 & 68.10 & 71.25 & 70.31 \\
                                   & VCR & \textbf{99.87} & 94.53 & 88.67 & 82.13 & 82.37 & 83.37 & 88.47 \\
\hline
\multirow{2}{*}{InternLM2.5-20B} & F1 & \textbf{63.52} & - & - & - & - & - & - \\
                                 & VCR & \textbf{67.27} & - & - & - & - & - & - \\
\hline
\multirow{2}{*}{Qwen2-7B} & F1 & \textbf{71.01} & 43.65 & 59.91 & 60.67 & 67.83 & 69.25 & 65.34 \\
                              & VCR & \textbf{99.97} & 99.70 & 99.00 & 99.77 & 98.87 & 99.10 & 98.30 \\
\hline
\multirow{2}{*}{Qwen2-72B} & F1 & 75.72 & 77.22 & \textbf{77.40} & - & - & - & - \\
                               & VCR & 97.73 & 98.77 & \textbf{99.47} & - & - & - & - \\
\hline
\multirow{2}{*}{Qwen2.5-3B} & F1 & 55.06 & \textbf{64.16} & 64.03 & 52.65 & 53.11 & 62.15 & 64.52 \\
                            & VCR & \textbf{100.00} & 99.73 & 99.50 & 93.17 & 90.40 & 71.73 & 57.40 \\
\hline
\multirow{2}{*}{Qwen2.5-7B} & F1 & 61.01 & 62.19 & 67.91 & 72.78 & \textbf{75.28} & 74.72 & 74.14 \\
                            & VCR & \textbf{100.00} & 99.80 & 97.97 & 80.70 & 80.50 & 85.47 & 75.63 \\
\hline
\multirow{2}{*}{Qwen2.5-72B} & F1 & 76.99 & 80.40 & \textbf{80.88} & - & - & - & - \\
                             & VCR & \textbf{99.40} & 94.17 & 91.23 & - & - & - & - \\
\hline
\end{tabular}
}
\label{tab:performance}
\end{table*}

\subsection{Fine-tuning Open-source LLMs}
We fine-tuned open-source LLMs using LoRA for smaller models (3B, 7B, and 8B) and QLoRA for larger models (20B, 70B, and 72B). The results, presented in Table~\ref{tab:finetuning_performance}, reveal several important findings:

\subsubsection{Dramatic Enhancement in Valid Classification Ratio}
Fine-tuning significantly improved the Valid Classification Ratio (VCR) across all models. Remarkably, after just 0.2 epochs, nearly all models achieved and maintained a perfect VCR score of 100\%. This rapid improvement demonstrates that fine-tuning swiftly enhances the models' ability to produce valid classifications, effectively addressing a key limitation observed in few-shot prompting approaches.

\subsubsection{F1 Score Progression and Mitigating Overfitting}
Our initial experiments involved fine-tuning the LLMs for 6 epochs, with checkpoints saved at the conclusion of each epoch. We observed a common pattern: F1 scores for most models initially increased and reached a peak, but tended to decline after 1 epoch. This trend suggested potential overfitting, where models began to memorize training data at the expense of generalization.

To address this challenge, we refined our approach in subsequent experiments by implementing a more granular checkpoint strategy, saving model states every 0.2 epochs. This adjustment allowed us to:

\begin{itemize}
    \item Capture more detailed performance data throughout the training process.
    \item Identify and select the optimal results with greater precision.
    \item Effectively mitigate the overfitting observed in the later stages of training.
\end{itemize}

This refined approach yielded improved results. For instance, Qwen2-72B achieved its peak F1 score of 80.33\% at 0.8 epochs, demonstrating the benefits of our adjusted fine-tuning strategy in optimizing model performance while minimizing overfitting.

\subsubsection{Superior Performance of Fine-tuned Open-source LLMs}

Fine-tuning has enabled several open-source LLMs to achieve F1 scores that not only match but also surpass the best few-shot performance of GPT-4o, which previously stood at 80.36\%. Among the most notable results:

\begin{itemize}
    \item \textbf{Llama3.1-70B} surpassed GPT-4o with a robust F1 score of 80.77\% after 1.0 epoch.
    \item \textbf{Qwen2-72B} delivered an outstanding F1 score of 80.33\% after 0.8 epochs.
    \item \textbf{InternLM2.5-20B} and \textbf{InternLM2.5-7B-1M} both attained impressive F1 scores of 80.28\% after 0.8 epochs.
\end{itemize}

% These results highlight that fine-tuned open-source models can not only compete with but also outperform proprietary models in complex reasoning tasks.

% \subsubsection{Efficiency of Fine-tuning}
% The rapid improvement in both VCR and F1 scores within the first epoch highlights the efficiency of fine-tuning for enhancing model performance on specific tasks. This efficiency is particularly valuable given the computational constraints observed in few-shot evaluations of larger models.

% \subsubsection{Balancing Performance and Generalization}
% The observed pattern of increasing then decreasing F1 scores during fine-tuning emphasizes the importance of careful monitoring and potential early stopping to prevent overfitting. This underscores the need for robust fine-tuning strategies that can balance performance improvements with generalization capabilities.

These findings collectively demonstrate the power of fine-tuning in unlocking the full potential of open-source LLMs, enabling them to compete with and even surpass proprietary models in challenging reasoning tasks. 

\begin{table*}[htbp]
\centering
\small
\caption{Fine-tuning Performance of Open-Source Language Models Over Training Epochs (F1 Score and Valid Classification Ratio in \%). \small{The top metrics for each model across different epochs are highlighted in bold. Fine-tuning led to significant improvements, with most models achieving a perfect VCR of 100\% within the first 0.2 epochs. Llama3.1-70B outperformed GPT-4o with an F1 score of 80.77\%, followed closely by Qwen2-72B, InternLM2.5-7B-1M, and InternLM2.5-20B, all of which closely matched GPT-4o's performance.}}

\resizebox{0.9\textwidth}{!}{
\begin{tabular}{l|l|ccccccccccc}
\hline
Model & Metric & 0.0 & 0.2 & 0.4 & 0.6 & 0.8 & 1.0 & 1.2 & 1.4 & 1.6 & 1.8 & 2.0 \\
\hline
\multirow{2}{*}{Llama3.1-8B} & F1 & 73.71 & 74.28 & 74.89 & 73.71 & 78.64 & \textbf{79.25} & 76.29 & 77.45 & 76.62 & 75.80 & 76.17 \\
                             & VCR & 80.33 & \textbf{100.00} & 99.93 & \textbf{100.00} & \textbf{100.00} & \textbf{100.00} & \textbf{100.00} & \textbf{100.00} & \textbf{100.00} & \textbf{100.00} & \textbf{100.00} \\
\hline
\multirow{2}{*}{Llama3.1-70B} & F1 & 75.26 & 79.11 & 76.25 & 75.61 & 78.15 & \textbf{80.77} & 76.44 & 78.20 & 78.34 & 78.24 & 77.92 \\
                              & VCR & 0.97 & \textbf{100.00} & \textbf{100.00} & \textbf{100.00} & \textbf{100.00} & \textbf{100.00} & \textbf{100.00} & \textbf{100.00} & \textbf{100.00} & \textbf{100.00} & \textbf{100.00} \\
\hline
\multirow{2}{*}{Mistral-7B} & F1 & 66.34 & 73.43 & \textbf{75.37} & 70.68 & 74.05 & 74.72 & 73.49 & 76.48 & 74.10 & 72.57 & 73.08 \\
                            & VCR & 1.17 & \textbf{100.00} & \textbf{100.00} & \textbf{100.00} & \textbf{100.00} & \textbf{100.00} & \textbf{100.00} & \textbf{100.00} & \textbf{100.00} & \textbf{100.00} & \textbf{100.00} \\
\hline
\multirow{2}{*}{InternLM2.5-7B} & F1 & 69.06 & 73.30 & 73.97 & 71.04 & 76.60 & 75.14 & 74.89 & \textbf{76.53} & 74.32 & 73.47 & 74.08 \\
                                & VCR & \textbf{100.00} & \textbf{100.00} & \textbf{100.00} & \textbf{100.00} & \textbf{100.00} & \textbf{100.00} & \textbf{100.00} & \textbf{100.00} & \textbf{100.00} & \textbf{100.00} & \textbf{100.00} \\
\hline
\multirow{2}{*}{InternLM2.5-7B-1M} & F1 & 52.45 & 78.65 & 78.87 & 75.66 & \textbf{80.28} & 78.43 & 78.77 & 77.77 & 77.56 & 77.40 & 77.78 \\
                                   & VCR & 99.87 & \textbf{100.00} & \textbf{100.00} & \textbf{100.00} & \textbf{100.00} & \textbf{100.00} & \textbf{100.00} & \textbf{100.00} & \textbf{100.00} & \textbf{100.00} & \textbf{100.00} \\
\hline
\multirow{2}{*}{InternLM2.5-20B} & F1 & 63.52 & 76.93 & 78.49 & 74.22 & \textbf{80.28} & 79.38 & 77.87 & 79.65 & 79.00 & 77.95 & 78.27 \\
                                 & VCR & 67.27 & \textbf{100.00} & \textbf{100.00} & \textbf{100.00} & \textbf{100.00} & \textbf{100.00} & \textbf{100.00} & \textbf{100.00} & \textbf{100.00} & \textbf{100.00} & \textbf{100.00} \\
\hline
\multirow{2}{*}{Qwen2-7B} & F1 & 71.01 & 74.90 & 77.49 & 73.32 & 74.76 & \textbf{78.89} & 73.63 & 75.54 & 72.92 & 71.64 & 72.60 \\
                            & VCR & \textbf{100.00} & 99.97 & \textbf{100.00} & \textbf{100.00} & \textbf{100.00} & 99.93 & \textbf{100.00} & 99.97 & \textbf{100.00} & \textbf{100.00} & \textbf{100.00} \\
\hline
\multirow{2}{*}{Qwen2-72B} & F1 & 75.72 & 78.28 & 77.01 & 78.40 & \textbf{80.33} & 77.00 & 79.01 & 79.06 & 79.30 & 80.42 & 79.74 \\
                             & VCR & 97.73 & \textbf{100.00} & \textbf{100.00} & \textbf{100.00} & \textbf{100.00} & \textbf{100.00} & \textbf{100.00} & \textbf{100.00} & \textbf{100.00} & \textbf{100.00} & \textbf{100.00} \\
\hline
\multirow{2}{*}{Qwen2.5-3B} & F1 & 53.43 & 69.74 & 69.97 & 71.77 & 72.64 & 72.33 & 71.73 & \textbf{74.10} & 73.56 & 72.63 & 72.90 \\
                            & VCR & \textbf{100.00} & 99.90 & \textbf{100.00} & \textbf{100.00} & \textbf{100.00} & \textbf{100.00} & \textbf{100.00} & \textbf{100.00} & \textbf{100.00} & \textbf{100.00} & \textbf{100.00} \\
\hline
\multirow{2}{*}{Qwen2.5-7B} & F1 & 61.01 & 74.83 & 75.77 & 76.94 & 76.66 & \textbf{78.15} & 76.61 & 76.55 & 76.89 & 76.81 & 77.04 \\
                            & VCR & \textbf{100.00} & 99.80 & 99.97 & 99.97 & \textbf{100.00} & 99.93 & \textbf{100.00} & 99.93 & 99.97 & 99.93 & 99.97 \\
\hline
\multirow{2}{*}{Qwen2.5-72B} & F1 & 76.99 & 78.74 & 77.19 & 77.83 & \textbf{79.84} & 76.82 & 79.01 & 78.26 & 78.24 & 78.94 & 78.50 \\
                             & VCR & 99.40 & \textbf{100.00} & \textbf{100.00} & \textbf{100.00} & \textbf{100.00} & \textbf{100.00} & \textbf{100.00} & \textbf{100.00} & \textbf{100.00} & \textbf{100.00} & \textbf{100.00} \\
\hline
\end{tabular}
}
\label{tab:finetuning_performance}
\end{table*}

\subsection{Comparative Analysis of Proprietary and Open-source LLMs}

Table \ref{tab:comparing_performance} presents a comprehensive comparison of large language models (LLMs) across various evaluation settings: baseline (zero-shot), best few-shot, and best fine-tuned configurations. The "best" refers to the highest F1 score achieved in each category. The table is structured in two sections: the first lists OpenAI models without fine-tuned results, while the second features open-source models with fine-tuning data. This organization facilitates a clear comparison between proprietary and open-source models across different evaluation strategies.

\subsubsection{Fine-tuning vs. Few-shot Learning}

The results indicate that task-specific fine-tuning generally enhances model performance, especially for open-source models. Most open-source models demonstrated better performance after fine-tuning compared to their few-shot results. However, the Qwen2.5-72B model stands out as an exception, achieving the highest overall performance with an F1 score of 80.88\% via few-shot learning, surpassing GPT-4o's best few-shot score of 80.36\%. Llama3.1-70B secured the second-highest overall performance with an F1 score of 80.77\% through fine-tuning, also surpassing GPT-4o.

These findings suggest that while fine-tuning is generally advantageous, few-shot learning can sometimes lead to superior outcomes. Researchers and practitioners should carefully evaluate both strategies when optimizing model performance for specialized tasks.

\subsubsection{Extended Context Windows}

The performance of the InternLM2.5-7B-1M model underscores the significant potential of extended context windows, especially in models with relatively fewer parameters. Despite its modest baseline performance (F1 score of 52.45\%), the model showed remarkable improvement in advanced settings. In the few-shot configuration, InternLM2.5-7B-1M achieved an F1 score of 77.01\%, ranking third among all evaluated open-source LLMs. After fine-tuning, the model further improved to an F1 score of 80.28\%, making it the third-highest performer overall in the fine-tuned category.

These results suggest that extended context windows can significantly enhance model performance, particularly when combined with fine-tuning. This finding points to a promising direction for future model development, where leveraging extended context windows might lead to improved reasoning capabilities without necessitating a large increase in parameter count. Future research should explore how to best utilize extended context windows in model architectures to further advance LLMs' capacity for complex reasoning.


\begin{table*}[htbp]
\centering
\caption{Comparative Analysis of Large Language Models (LLMs) across Baseline (Zero-shot), Best Few-shot, and Best Fine-tuned Configurations (F1 Score and Valid Classification Ratio in \%). \small{"Best" refers to the highest F1 score achieved in each evaluation setting. The table is divided into two sections: OpenAI models (top) and open-source models (bottom). For each section, the highest metrics in each category are highlighted in bold. Notable observations include Qwen2.5-72B's superior few-shot performance with an F1 score of 80.88\%, surpassing GPT-4o's best score of 80.36\%. Among fine-tuned models, Llama3.1-70B achieves the highest F1 score of 80.77\%. The InternLM2.5-7B-1M model demonstrates significant improvement from a baseline F1 of 52.45\% to 77.01\% in few-shot and 80.28\% after fine-tuning, highlighting the potential of extended context windows.}}
\label{tab:comparing_performance}
\resizebox{0.6\textwidth}{!}{
\begin{tabular}{lcccccccccc}
\hline
\multirow{2}{*}{Model} & \multicolumn{2}{c}{Baseline (\%)} & \multicolumn{4}{c}{Best Few-shot (\%)} & \multicolumn{4}{c}{Best Fine-tuned (\%)} \\
\cmidrule(lr){2-3} \cmidrule(lr){4-7} \cmidrule(lr){8-11}
& F1 & VCR & F1 & VCR & Shots & & F1 & VCR & Epoch & \\
\hline
GPT-4o-mini & 72.96 & 99.17 & 72.96 & 99.17 & 0 & & - & - & - & \\
GPT-4o & \textbf{79.53} & 6.60 & \textbf{80.36} & \textbf{99.97} & 10 & & - & - & - & \\
o1-mini & 73.77 & \textbf{99.90} & 75.90 & 99.77 & 50 & & - & - & - & \\
o1-preview & 74.51 & 99.80 & 77.08 & 98.17 & 50 & & - & - & - & \\
\hline
Llama3.1-8B & 73.71 & 80.33 & 76.97 & 73.27 & 30 & & 79.25 & \textbf{100.00} & 1.0 & \\
Llama3.1-70B & 75.26 & 0.97 & 75.67 & 83.27 & 10 & & \textbf{80.77} & \textbf{100.00} & 1.0 & \\
Mistral-7B & 66.34 & 1.17 & 68.63 & 7.00 & 30 & & 76.48 & \textbf{100.00} & 1.4 & \\
InternLM2.5-7B & 69.06 & \textbf{100.00} & 72.71 & 99.90 & 5 & & 76.60 & \textbf{100.00} & 0.8 & \\
InternLM2.5-7B-1M & 52.45 & 99.87 & 77.01 & 94.53 & 5 & & 80.28 & \textbf{100.00} & 0.8 & \\
InternLM2.5-20B & 63.52 & 67.27 & - & - & - & & 80.28 & \textbf{100.00} & 0.8 & \\
Qwen2-7B & 71.01 & 99.97 & 71.01 & 99.97 & 0 & & 77.49 & \textbf{100.00} & 0.4 & \\
Qwen2-72B & 75.72 & 97.73 & 77.40 & 99.47 & 10 & & 80.33 & \textbf{100.00} & 0.8 & \\
Qwen2.5-3B & 55.06 & \textbf{100.00} & 64.16 & 99.73 & 5 & & 74.10 & \textbf{100.00} & 1.4 & \\
Qwen2.5-7B & 61.01 & \textbf{100.00} & 75.28 & 80.50 & 30 & & 78.15 & 99.93 & 1.0 & \\
Qwen2.5-72B & \textbf{76.99} & 99.40 & \textbf{80.88} & 91.23 & 10 & & 79.84 & \textbf{100.00} & 0.8 & \\
\hline
\end{tabular}
}
\end{table*}



\section{Conclusion}

This study presents a comprehensive evaluation of large language models (LLMs) in tackling the Turtle Soup Question Answering (TSQA) task, highlighting both the strengths and limitations of different approaches, including few-shot prompting and fine-tuning. Our findings challenge the notion that proprietary models are inherently superior, demonstrating that fine-tuned open-source models can achieve, and in some cases, exceed the performance of established models like GPT-4o.

Key insights from our research include:

\begin{itemize}
    \item \textbf{Advantages of Fine-tuning:} Fine-tuning significantly enhances model performance, particularly in terms of Valid Classification Ratio (VCR), thereby increasing the reliability of models in producing valid outputs. Most open-source models exhibited substantial improvements after fine-tuning, often surpassing their few-shot performance.
    
    \item \textbf{Potential of Few-shot Learning:} Although fine-tuning generally led to better results, few-shot learning also proved highly effective. For example, Qwen2.5-72B achieved the highest F1 score of 80.88\%, outperforming GPT-4o. This demonstrates that few-shot learning remains a viable approach, especially when fine-tuning is not feasible.
    
    \item \textbf{Impact of Extended Context Windows:} The strong performance of InternLM2.5-7B-1M, equipped with an extended context length of 1 million tokens, underscores the potential of this feature. The model showed significant improvement from a modest baseline, highlighting the importance of extended context windows in enhancing logical reasoning capabilities.
    
    \item \textbf{Balancing Optimization and Generalization:} The study emphasizes the need to optimize fine-tuning strategies to balance performance gains with generalization, as overfitting was observed when models were trained beyond 1 epoch.
\end{itemize}

These findings underscore the potential of fine-tuned open-source LLMs to perform complex reasoning tasks at levels comparable to, or even surpassing, proprietary models.

Future research should focus on refining fine-tuning techniques to maximize performance while mitigating overfitting. Additionally, exploring the impact of extended context windows across a broader range of models and tasks could offer further insights into leveraging this feature to enhance reasoning capabilities. Developing more efficient and scalable evaluation methods will be crucial as models continue to grow in size and complexity. The integration of few-shot learning with fine-tuning might also provide a hybrid approach that captures the benefits of both strategies, especially in resource-constrained environments.

In conclusion, while substantial progress has been made in enhancing the reasoning capabilities of LLMs, there remains significant potential for further development. Future efforts should focus on refining training methodologies to improve efficiency, optimizing model architectures specifically for complex reasoning tasks, and establishing robust evaluation frameworks that comprehensively assess model performance. By addressing these challenges, we can continue to advance the field of AI, moving closer to developing models that not only rival but potentially exceed human-level reasoning capabilities across a diverse range of problem-solving scenarios.


% \section*{Acknowledgment}

\bibliographystyle{IEEEtran}
\bibliography{references}

\end{document}
